\chapter{Introduction}

\section{Introduction}

The fast spread of Internet of Things (IoT) technologies has transformed modern digital ecosystems by enabling seamless connectivity among billions of heterogeneous devices. IoT systems are increasingly deployed across diverse application domains, including smart cities, healthcare, industrial automation, transportation, energy management, and critical infrastructure. These systems generate enormous amounts of data and facilitate real-time decision-making, significantly improving operational efficiency and service delivery.


Despite these advantages, the widespread adoption of IoT introduces substantial cybersecurity challenges. Traditional security models, which rely on perimeter-based protection and implicit trust assumptions, are becoming ineffective in highly distributed and dynamic environments. IoT devices frequently interact across multiple administrative and security domains, creating complex trust relationships and expanding the attack surface. As a result, cyber threats such as unauthorized access, lateral movement attacks, identity spoofing, insider threats, and session hijacking have become increasingly prevalent.


To address these challenges, Zero Trust Architecture (ZTA) has emerged as a modern security paradigm founded on the principle of “Never Trust, Always Verify.” Unlike conventional security approaches, ZTA continuously verifies users, devices, and applications before granting access to resources. While ZTA significantly improves security, implementing it in large-scale cross-domain IoT environments is still challenging due to frequent authentication requirements, trust management complexity, scalability concerns, and communication overhead.


Recent advances in Federated Learning (FL) offer promising opportunities for addressing these challenges. FL enables multiple entities to collaboratively train machine learning models without sharing raw data, thereby preserving privacy while supporting distributed intelligence. By integrating FL with Zero Trust principles, adaptive trust management mechanisms can be developed to facilitate efficient cross-domain authentication, dynamic trust evaluation, and scalable authorization decisions.


This research presents a Federated Learning-Based Adaptive Trust Management Framework for Cross-Domain Zero Trust IoT Networks to enhance security, privacy, scalability, and trust coordination across heterogeneous IoT environments.


\section{Background of the Study}

Cybersecurity has moved from perimeter-based security models to identity-centric security architectures. Older methods of network security were based on the idea that all entities inside a trusted network boundary could be considered trustworthy. But the development of cloud computing, remote work, mobile devices, and IoT ecosystems has made that assumption no longer valid.


Zero Trust Architecture, as formally defined by the National Institute of Standards and Technology (NIST SP 800-207), gets rid of any assumed trust and calls for the continuous verification of each and every access request regardless of where the request comes from. This approach includes principles such as least-privilege access, continuous authentication, dynamic authorization, micro-segmentation, and continuous monitoring.


In IoT environments, devices commonly communicate over a number of domains, for example, enterprise networks, cloud platforms, industrial systems, and service provider infrastructures. Currently, the cross-domain authentication methods in use usually involve full re-authentication each time a device enters a new security domain. This process introduces significant latency and computational overhead while limiting interoperability.


Trust management has become as a critical component of ZTA, enabling security systems to assess the trustworthiness of devices and users dynamically. However, most existing trust management approaches depend on centralized architectures that face scalability limitations and privacy concerns.


Federated Learning provides a decentralized alternative by enabling collaborative trust evaluation without exposing sensitive device data. Through distributed model training and secure parameter exchange, FL can support adaptive trust assessment while preserving privacy. The integration of FL with ZTA therefore represents a promising direction for improving cross-domain trust management in large-scale IoT networks.


\section{Problem Statement}

The growing complexity of cross-domain IoT ecosystems presents significant challenges for existing Zero Trust implementations. Current authentication and authorization mechanisms often suffer from excessive communication overhead, repeated authentication processes, high latency, and limited trust portability across security domains.


Furthermore, existing trust management systems are predominantly centralized, which makes them vulnerable to single points of failure, scalability limitations, and privacy risks. While Federated Learning has demonstrated effectiveness in privacy-preserving collaborative learning, its application to adaptive trust management within cross-domain Zero Trust IoT environments remains insufficiently explored.


Consequently, there is a need for a scalable, privacy-preserving, and adaptive trust management framework capable of dynamically evaluating trust across multiple security domains while supporting efficient authentication and authorization decisions in Zero Trust IoT networks.


\section{Significance of the Study}

This study is significant for several reasons:


It contributes to the improvement of Zero Trust security frameworks for large-scale IoT environments.


It addresses the missing research on cross-domain trust management in current literature.


It leverages Federated Learning to enhance privacy preservation by eliminating the need to exchange raw device data.


It improves authentication efficiency by enabling adaptive trust evaluation and trust migration across security domains.


It provides a scalable framework suitable for heterogeneous IoT deployments involving multiple administrative domains.


It contributes to both academic research and practical cybersecurity implementations in smart cities, healthcare systems, industrial IoT, and critical infrastructure.


\section{Objectives of the Study}

\paragraph{General Objective}

To develop a Federated Learning-Based Adaptive Trust Management Framework that enhances cross-domain authentication and authorization in Zero Trust IoT networks.


\paragraph{Specific Objectives}

To analyze existing trust management approaches in Zero Trust IoT environments.


To design a federated trust management framework capable of supporting cross-domain trust evaluation in a Zero Trust IoT environment.


To develop an adaptive trust scoring mechanism based on IoT device behavior and contextual information.


To integrate Federated Learning into the trust evaluation process while preserving privacy.


To evaluate the framework in terms of security, scalability, authentication latency, trust accuracy, and communication overhead.


To compare the proposed framework with existing Zero Trust and trust management solutions.


\section{Research Questions}

The study seeks to answer the following research questions:


RQ1: What are the limitations of existing trust management approaches in cross-domain Zero Trust IoT environments?


RQ2: How can Federated Learning be utilized to facilitate privacy-preserving trust evaluation across multiple security domains?


RQ3: How can adaptive trust scores be generated to support dynamic authentication and authorization decisions?


RQ4: To what extent can the proposed framework reduce authentication latency and communication overhead?


RQ5: How does the proposed framework compare to other existing Zero Trust trust management solutions in terms of security, scalability, and efficiency?


\section{Limitations of the Study}

This study is subject to the following limitations:


The proposed framework will be evaluated primarily through simulation and experimental environments rather than large-scale real-world deployment.


The research focuses on IoT networks operating under Zero Trust principles and may not directly address other network architectures.


The Federated Learning model assumes participating domains are cooperative and willing to exchange model parameters.


The study does not address post-quantum cryptographic requirements.


Resource-constrained IoT devices may require additional optimization mechanisms beyond the scope of this research.


\section{Methodology of the Study}

This study adopts a Design Science Research methodology combined with quantitative experimental evaluation.


The research process consists of the following phases:


Literature Review: Conduct a systematic review of Zero Trust Architecture, trust management, Federated Learning, and cross-domain authentication.


Requirement Analysis: Identify security, privacy, scalability, and trust management requirements.


Framework Design: Design the proposed federated trust management architecture and trust evaluation mechanisms.


Model Development: Develop adaptive trust scoring algorithms and federated learning processes.


Simulation and Implementation: Implement the framework using simulation tools and machine learning environments.


Performance Evaluation: Evaluate the framework using metrics such as authentication latency, trust accuracy, communication overhead, throughput, scalability, and security resilience.


Comparative Analysis: Compare the proposed framework against existing Zero Trust and trust management approaches.


\section{Thesis Structure}

This thesis is organized into six chapters as follows:


Chapter One presents the introduction, background, problem statement, significance, objectives, research questions, methodology, limitations, and overall thesis organization.


Chapter Two reviews the relevant literature on Zero Trust Architecture, IoT security, trust management systems, Federated Learning, and cross-domain authentication mechanisms.


Chapter Three presents the conceptual framework, system architecture, assumptions, and design requirements of the proposed solution.


Chapter Four describes the proposed Federated Learning-Based Adaptive Trust Management Framework, including trust evaluation algorithms and cross-domain trust management mechanisms.


Chapter Five presents the implementation details, experimental setup, simulation results, performance evaluation, and comparative analysis.


Chapter Six summarizes the research findings, contributions, limitations, and recommendations for future work.
